\documentclass[11pt,a4paper]{article}
\usepackage[utf8]{inputenc}
\usepackage[T1]{fontenc}
\usepackage[ngerman]{babel}
\usepackage{amsmath,amssymb,graphicx,booktabs,xcolor}
\usepackage[margin=25mm]{geometry}
\usepackage{tcolorbox,hyperref}

\definecolor{darkblue}{HTML}{16213e}
\definecolor{keygreen}{HTML}{2e7d32}

\title{\textcolor{darkblue}{Dok. 0018\\Hörbarkeit der Biegemoden:\\Kammer-Saugkreis, Transiente, Torsion}}
\author{Harmonikabau --- Technische Dokumentation}
\date{}

\begin{document}
\maketitle

\noindent\textit{Korrektur zu Dok.~0016: Die Biegemoden können unter bestimmten Bedingungen
hörbar werden. Drei Mechanismen: Kammer-Resonanz, Einschwingvorgang, Torsion.
Referenz: Dok.~0009, 0011, 0015--0017.}

\noindent\textit{\small Berechnungsskript:
\textbf{pruefskript\_0018\_hoerbarkeit.py} (modale Partizipation,
Kammer-Saugkreis, Transiente, Torsionsmoden) im selben Verzeichnis auf GitHub.}

\section{Korrektur: Biegemoden sind nicht immer unhörbar}

Dok.~0016 hat gezeigt, dass die inharmonischen Biegemoden im Normalfall
30--50\,dB leiser sind als die harmonischen Obertöne.
Die Praxis zeigt aber: Mit einer abstimmbaren Resonanzkammer kann man
Mode~2 bei manchen Zungen so stark anregen, dass \textbf{nur Mode~2 hörbar ist}.
Und beim Einschwingvorgang prägt Mode~2 den Klangcharakter.

\begin{tcolorbox}[colframe=red!60!black, colback=red!5!white]
\textbf{Die Grundaussage bleibt:} Die harmonischen Obertöne dominieren
im eingeschwungenen Zustand. Aber Mode~2 ist nicht vernachlässigbar.
\end{tcolorbox}

\section{Kammer als Saugkreis: Mode 2 wird hörbar durch Subtraktion}

Die Kammer wirkt als \textbf{Saugkreis} (Kerbfilter, Dok.~0005--0008) ---
sie \textbf{unterdrückt} Frequenzen bei $f_H$. Die Kammer verstärkt Mode~2 nicht.
Mode~2 wird hörbar, weil der Grundton und seine Harmonischen gedämpft werden,
während Mode~2 \textbf{außerhalb der Kerbe} liegt.

Voraussetzung: (a) $f_H$ trifft einen Harmonischen von $f_1$ $\to$ Grundton
wird gedämpft. (b) $f_2$ (Biegemode) liegt \textbf{nicht} in der Kerbe
$\to$ Mode~2 bleibt übrig. Das ist Subtraktion, nicht Addition.

\textbf{Warum geht das nicht bei jeder Zunge?} Die Profilierung bestimmt $f_2$.
Bei gleichmäßiger Zunge liegt $f_2 \approx 6{,}27 \times f_1$ --- nahe an
$6 \times f_1$ (Harmonischer). Bei stärkerer Profilierung ($f_2 \approx 4$--$5 \times f_1$)
liegt Mode~2 zwischen den Harmonischen und überlebt die Kerbe.

\begin{tcolorbox}[colframe=keygreen, colback=green!5!white]
\textbf{Im Akkordeon trifft es die Töne mit Ansprache-Problemen:}
Dort ist die Kammer-Kopplung am stärksten, die Stimmplatten sind stärker
profiliert, und das Zusammenspiel von Profilierung und Kammergeometrie
erfüllt zufällig die Bedingung $f_H \neq f_2$.
\end{tcolorbox}

\section{Transiente Anregung: Einschwingvorgang}

Beim Drücken einer Taste trifft der Balgdruck als Sprungfunktion.
Jede Biegemode wird initial angestoßen.
Mode~2 startet bei $\approx -25$\,dB und klingt in $\approx 40$\,ms ab.

\begin{tcolorbox}[colframe=keygreen, colback=green!5!white]
\textbf{Mode~2 ist im Einschwingvorgang hörbar} --- als kurze Klangfarbe
in den ersten 50--100\,ms. Verschiedene Zungen haben verschiedene $f_2$
$\to$ verschiedener Klangcharakter.
\end{tcolorbox}

\section{Torsionsschwingungen}

Die Torsionsmoden liegen höher als die Biegemoden: $f_{T1} \approx 14$--$22 \times f_1$.
Die Torsion selbst ist als Ton kaum hörbar --- aber die seitliche Auslenkung
führt bei Basszungen zu Kanalberührung.

Bei 1° Torsion: Basszunge (W=8\,mm): $\Delta y = 0{,}07$\,mm $>$ Spalt 0,04\,mm
$\to$ Kontakt möglich. \textbf{Im normalen Betrieb sollte Kratzen nicht vorkommen} ---
die Spalte sind groß genug. Aber bei \textbf{kalten Instrumenten} ist Kratzen normal:
Das Metall kontrahiert ($\approx 2$--$3$\,µm bei $\Delta T = -17$\,K),
der Spalt wird enger, die Toleranz sinkt. Sobald das Instrument sich erwärmt,
verschwindet das Kratzen.

\begin{tcolorbox}[colframe=keygreen, colback=green!5!white]
\textbf{Paradoxon:} Die passgenauesten Stimmplatten --- also die teuersten ---
sind am anfälligsten für Kratzen bei Kälte. Enge Spalte (bessere Ansprache)
bedeuten weniger Toleranz für Torsion. Eine Platte mit 25\,µm Spalt verliert
bei 3\,µm Kontraktion 12\,\% Spielraum --- eine mit 50\,µm nur 6\,\%.
Die billige Platte mit großem Spalt kratzt nie, die teure mit perfektem Spalt
kratzt im Winter.
\end{tcolorbox}

\section{Zusammenfassung}

\begin{enumerate}
\item \textbf{Eingeschwungen, ohne Kammer-Treffer:} Mode~2 bei $-37$\,dB ---
  nicht hörbar. Harmonische dominieren.
\item \textbf{Kammer als Saugkreis:} Die Kammer dämpft Harmonische.
  Wenn $f_2$ außerhalb der Kerbe liegt, bleibt Mode~2 übrig --- hörbar
  durch Subtraktion. Trifft besonders Töne mit Ansprache-Problemen.
\item \textbf{Transiente:} Mode~2 bei $-25$\,dB, Dauer $\approx 40$\,ms.
  Prägt Klangcharakter.
\item \textbf{Torsion:} $f_{T1} \approx 14$--$22 \times f_1$.
  Kanalberührung im Normalbetrieb nicht vorgesehen.
  Bei \textbf{kalten Instrumenten} normal --- Metall kontrahiert, Spalt enger.
  Verschwindet beim Erwärmen.
\item \textbf{Korrektur:} ``Biegemoden nicht hörbar'' gilt nur eingeschwungen
  ohne Kammer-Treffer.
\end{enumerate}

\bigskip
\textit{Die Biegemoden sind leise --- aber nicht stumm.
Die Kammer kann sie wecken, der Transient lässt sie kurz erklingen,
die Torsion macht sie als Kratzen hörbar.}

\end{document}
